\nonstopmode{}
\documentclass[a4paper]{book}
\usepackage[times,inconsolata,hyper]{Rd}
\usepackage{makeidx}
\makeatletter\@ifl@t@r\fmtversion{2018/04/01}{}{\usepackage[utf8]{inputenc}}\makeatother
% \usepackage{graphicx} % @USE GRAPHICX@
\makeindex{}
\begin{document}
\chapter*{}
\begin{center}
{\textbf{\huge Package `msarima'}}
\par\bigskip{\large \today}
\end{center}
\ifthenelse{\boolean{Rd@use@hyper}}{\hypersetup{pdftitle = {msarima: Multi-Seasonal ARIMA Models}}}{}
\ifthenelse{\boolean{Rd@use@hyper}}{\hypersetup{pdfauthor = {Daniele Girolimetto; Francesco Lisi; Matteo Grigoletto}}}{}
\begin{description}
\raggedright{}
\item[Title]\AsIs{Multi-Seasonal ARIMA Models}
\item[Version]\AsIs{0.1.0}
\item[Description]\AsIs{Fit, forecast and simulate autoregressive integrated moving
average (ARIMA) models with more than one seasonal component, following
Lisi and Grigoletto (2025) <}\Rhref{https://doi.org/10.1002/for.3283}{doi:10.1002/for.3283}\AsIs{>. Series observed at high
frequency often display several seasonal cycles at once, such as the daily
and the weekly pattern of hourly data, which a single seasonal ARIMA
operator cannot represent. Estimation is performed by conditional maximum
likelihood through numerical optimisation, with support for fixed
coefficients and a constant term. Multi-seasonal Ljung-Box and Monti
portmanteau tests, following Grigoletto and Lisi (2026)
<}\Rhref{https://doi.org/10.48550/arXiv.2609.32381}{doi:10.48550/arXiv.2609.32381}\AsIs{>, seasonal correlograms, forecasts with
prediction intervals, and simulation from a model specification or from a
fitted model are also provided.}
\item[License]\AsIs{GPL-3}
\item[Encoding]\AsIs{UTF-8}
\item[Roxygen]\AsIs{list(markdown = TRUE)}
\item[RoxygenNote]\AsIs{7.3.3}
\item[Depends]\AsIs{R (>= 4.1.0),
forecast}
\item[Imports]\AsIs{polynom,
stats}
\item[Suggests]\AsIs{ggplot2,
knitr,
rmarkdown,
testthat (>= 3.0.0)}
\item[VignetteBuilder]\AsIs{knitr}
\item[Config/testthat/edition]\AsIs{3}
\item[URL]\AsIs{}\url{https://github.com/danigiro/msarima}\AsIs{, }\url{https://danigiro.github.io/msarima}\AsIs{}
\item[BugReports]\AsIs{}\url{https://github.com/danigiro/msarima/issues}\AsIs{}
\item[LazyData]\AsIs{true}
\end{description}
\Rdcontents{Contents}
\HeaderA{msarima-package}{msarima: Multi-Seasonal ARIMA Models}{msarima.Rdash.package}
\keyword{package}{msarima-package}
%
\begin{Description}
Fit, forecast and simulate autoregressive integrated moving average (ARIMA) models with more than one seasonal component, following Lisi and Grigoletto (2025) \Rhref{https://doi.org/10.1002/for.3283}{doi:10.1002\slash{}for.3283}. Series observed at high frequency often display several seasonal cycles at once, such as the daily and the weekly pattern of hourly data, which a single seasonal ARIMA operator cannot represent. Estimation is performed by conditional maximum likelihood through numerical optimisation, with support for fixed coefficients and a constant term. Multi-seasonal Ljung-Box and Monti portmanteau tests, following Grigoletto and Lisi (2026) \Rhref{https://doi.org/10.48550/arXiv.2609.32381}{doi:10.48550\slash{}arXiv.2609.32381}, seasonal correlograms, forecasts with prediction intervals, and simulation from a model specification or from a fitted model are also provided.
\end{Description}
%
\begin{Author}
\strong{Maintainer}: Daniele Girolimetto \email{daniele.girolimetto@unipd.it} (\Rhref{https://orcid.org/0000-0001-9387-1232}{ORCID}) [copyright holder]

Authors:
\begin{itemize}

\item{} Francesco Lisi \email{francesco.lisi@unipd.it} (\Rhref{https://orcid.org/0000-0002-8749-8889}{ORCID})
\item{} Matteo Grigoletto \email{matteo.grigoletto@unipd.it} (\Rhref{https://orcid.org/0000-0001-9881-2039}{ORCID})

\end{itemize}


\end{Author}
%
\begin{SeeAlso}
Useful links:
\begin{itemize}

\item{} \url{https://github.com/danigiro/msarima}
\item{} \url{https://danigiro.github.io/msarima}
\item{} Report bugs at \url{https://github.com/danigiro/msarima/issues}

\end{itemize}


\end{SeeAlso}
\HeaderA{forecast.msarima}{Forecast from a Multi-Seasonal ARIMA Model}{forecast.msarima}
%
\begin{Description}
Returns forecasts and prediction intervals for a fitted multi-seasonal ARIMA model.
\end{Description}
%
\begin{Usage}
\begin{verbatim}
## S3 method for class 'msarima'
forecast(object, h = 10, level = c(80, 95), fan = FALSE, ...)
\end{verbatim}
\end{Usage}
%
\begin{Arguments}
\begin{ldescription}
\item[\code{object}] An object of class \code{"msarima"}.

\item[\code{h}] Number of periods to forecast. Default is 10.

\item[\code{level}] Confidence levels for the prediction intervals, in percent.
Default is \code{c(80, 95)}.

\item[\code{fan}] If \code{TRUE}, \code{level} is set to \code{seq(51, 99, by = 3)},
which is suitable for fan plots. Default is \code{FALSE}.

\item[\code{...}] Additional arguments (currently not used).
\end{ldescription}
\end{Arguments}
%
\begin{Details}
Point forecasts are computed by running the model recursion forward, setting the
future innovations to zero. Differencing and the constant term, when present, are
taken from the fitted model, so the forecasts are on the scale of the original
series.

Prediction intervals assume Gaussian innovations. The forecast variance
\eqn{V_h}{} is obtained from the MA(\eqn{\infty}{}) representation of the model.
\end{Details}
%
\begin{Value}
An object of class \code{"forecast"} containing the following elements:
\begin{ldescription}
\item[\code{method}] The name of the forecasting method (same as \code{object\$method}).
\item[\code{model}] The fitted \code{"msarima"} model object.
\item[\code{level}] The confidence levels for prediction intervals (as percentages).
\item[\code{mean}] Point forecasts as a time series (ts or msts object).
\item[\code{lower}] Lower limits for prediction intervals. Matrix with one column for each
confidence level specified in \code{level}.
\item[\code{upper}] Upper limits for prediction intervals. Matrix with one column for each
confidence level specified in \code{level}.
\item[\code{x}] The original time series.
\item[\code{series}] The name of the series.
\item[\code{fitted}] Fitted values from the model.
\item[\code{residuals}] Residuals from the model.
\end{ldescription}
\end{Value}
%
\begin{SeeAlso}
\code{\LinkA{msarima}{msarima}}, \code{\LinkA{simulate.msarima}{simulate.msarima}}
\end{SeeAlso}
%
\begin{Examples}
\begin{ExampleCode}
# Simulate and fit model
set.seed(123)
y <- msarima.sim(nsim = 200,
                 order = c(1, 0, 0),
                 seasonal = list(c(1, 0, 0), c(0, 0, 1)),
                 seasonal.period = c(4, 11),
                 ar = c(0.5, 0.5), ma = 0.3, sigma2 = 1)

fit <- msarima(y, order = c(1, 0, 0),
               seasonal = list(c(1, 0, 0), c(0, 0, 1)),
               seasonal.period = c(4, 11))

# Forecast 20 steps ahead with default confidence levels (80%, 95%)
fc <- forecast(fit, h = 20)
print(fc)

# Forecast with custom prediction intervals
fc2 <- forecast(fit, h = 20, level = c(90, 95, 99))

\end{ExampleCode}
\end{Examples}
\HeaderA{MAD\_m40}{Hourly Traffic Volume on the M-40 Ring Road, Madrid}{MAD.Rul.m40}
\keyword{datasets}{MAD\_m40}
%
\begin{Description}
Number of vehicles per hour recorded by a single traffic sensor on the M-40, the
outer ring road of Madrid, over the eight weeks from 3 February to 29 March 2024.

The series shows the two seasonal patterns that \code{\LinkA{msarima}{msarima}} is meant for:
a daily cycle of period 24, with the morning and evening commuting peaks, and a
weekly cycle of period 168, with markedly lower volumes at the weekend.
\end{Description}
%
\begin{Usage}
\begin{verbatim}
MAD_m40
\end{verbatim}
\end{Usage}
%
\begin{Format}
An \code{msts} object of length 1344 (56 days, or eight complete weeks, of hourly
data) with seasonal periods 24 and 168. The first observation is midnight of
Saturday 3 February 2024 and the last is 11 p.m. of Friday 29 March 2024.
\end{Format}
%
\begin{Details}
The raw measurements are half-hourly counts from sensor 3753, recorded between
1 February and 31 May 2024 and summed in pairs to obtain hourly volumes; the eight
weeks kept here are the window studied by Grigoletto and Lisi (2026), who chose this
sensor for having the fewest missing observations in the network. A few of the
underlying half-hourly counts had been interpolated by the provider, so 3 of the
hourly values are not whole numbers. There are no missing values and no gaps in the
hourly grid.

The script that builds the object from the raw file is available at
\url{https://github.com/danigiro/msarima/blob/main/data-raw/MAD_m40.R}.
\end{Details}
%
\begin{Source}
Madrid City Council open data portal, \url{https://datos.madrid.es}.
\end{Source}
%
\begin{References}
Grigoletto, M. and Lisi, F. (2026). Portmanteau goodness-of-fit tests in the
presence of simple, mixed, and multiple seasonal autocorrelation. \emph{arXiv}.
\Rhref{https://doi.org/10.48550/arXiv.2609.32381}{doi:10.48550\slash{}arXiv.2609.32381}
\end{References}
%
\begin{Examples}
\begin{ExampleCode}
# The daily and weekly cycles
plot(window(MAD_m40, end = 4), ylab = "vehicles per hour")

# A multi-seasonal ARIMA with one daily and one weekly difference

fit <- msarima(MAD_m40,
               order = c(1, 0, 2),
               seasonal = list(c(0, 1, 1), c(0, 1, 1)))
summary(fit)
plot(forecast(fit, h = 48))

\end{ExampleCode}
\end{Examples}
\HeaderA{msarima}{Fit a Multi-Seasonal ARIMA Model}{msarima}
%
\begin{Description}
Fits an ARIMA model with one or more seasonal components to a univariate time series.
This function extends the standard ARIMA framework to time series with several
seasonal patterns, such as daily and weekly seasonality in hourly data, or weekly
and yearly seasonality in daily data.
\end{Description}
%
\begin{Usage}
\begin{verbatim}
msarima(
  y,
  order = c(0, 0, 0),
  seasonal = list(),
  seasonal.period = NULL,
  fixed = NULL,
  include.constant = FALSE,
  optim_control = list(maxit = 200)
)
\end{verbatim}
\end{Usage}
%
\begin{Arguments}
\begin{ldescription}
\item[\code{y}] A numeric vector or time series object (\code{ts} or \code{msts}) containing
the data to be modelled. If \code{y} is a \code{ts} object and \code{seasonal.period}
is not supplied, its frequency is used as the first seasonal period; if \code{y} is
an \code{msts} object, its seasonal periods are used. Missing values are not
allowed.

\item[\code{order}] A numeric vector of length 3 specifying the non-seasonal part of the
ARIMA model: the AR order (p), the degree of differencing (d), and the MA order (q).
Default is \code{c(0, 0, 0)}.

\item[\code{seasonal}] A list of numeric vectors, each of length 3, specifying the seasonal
components. Each vector gives the seasonal AR order (P), the seasonal differencing
order (D), and the seasonal MA order (Q) of one component. For example,
\code{list(c(1, 0, 1), c(0, 1, 1))} specifies two seasonal components. Use
\code{list()} for a model without seasonal components. Default is \code{list()}.

\item[\code{seasonal.period}] A numeric vector of seasonal periods, one for each element of
\code{seasonal}, for example \code{c(7, 30)} for weekly and monthly seasonality.
If \code{NULL}, the periods are taken from \code{y} when it is a \code{ts} or
\code{msts} object. Periods can be given in any order but must be distinct.
Default is \code{NULL}.

\item[\code{fixed}] A numeric vector with one entry per parameter, used to hold some of them
at a given value: only the \code{NA} entries are estimated. The parameters are
ordered as non-seasonal AR, seasonal AR (one block per seasonal component),
non-seasonal MA, seasonal MA (one block per seasonal component), sigma, and the
constant term when \code{include.constant = TRUE}. Default is \code{NULL}, meaning
that all parameters are estimated.

\item[\code{include.constant}] Should the model include a constant term? The constant is the
intercept of the model equation, not the mean of the series: for a stationary model
the two are related by \eqn{\mu = c / \phi(1)}{}, where \eqn{c}{} is the intercept and
\eqn{\phi(1)}{} is the full AR polynomial (seasonal components included) evaluated
at 1. For example, for an AR(1) model \eqn{\mu = c / (1 - \phi_1)}{}. A constant
cannot be combined with differencing and is dropped, with a warning, when \code{d}
or any \code{D} is non-zero. Default is \code{FALSE}.

\item[\code{optim\_control}] A list of control parameters passed to \code{\LinkA{optim}{optim}}.
Default is \code{list(maxit = 200)}.
\end{ldescription}
\end{Arguments}
%
\begin{Details}
The model is estimated by conditional maximum likelihood: the residuals are computed
recursively and the resulting Gaussian log-likelihood is maximised numerically with
\code{\LinkA{optim}{optim}}, using the BFGS algorithm.

When \code{d} or any \code{D} is non-zero, the series is differenced before estimation
and the ARMA part is estimated on the differenced data. The first
\code{d + sum(D * S)} observations are therefore not used, which is why the residuals
and the fitted values start with as many \code{NA} values: these observations have
no differenced value to predict, so no residual is defined for them under the
conditional likelihood.

Standard errors are computed from the inverse of the observed information matrix
(negative Hessian). If the Hessian is singular, a warning is issued and the standard
errors are set to \code{NA}.
\end{Details}
%
\begin{Value}
An object of class \code{"msarima"}, which is a list containing:
\begin{ldescription}
\item[\code{order}] The non-seasonal ARIMA order (p, d, q).
\item[\code{seasonal}] List of seasonal orders (P, D, Q), one vector per seasonal component.
\item[\code{seasonal.period}] Vector of seasonal periods.
\item[\code{coef}] Named vector of estimated coefficients: the AR, SAR, MA and SMA
parameters, followed by the intercept if \code{include.constant = TRUE}.
\item[\code{var\_coef}] Variance-covariance matrix of the coefficient estimates, with
the same dimension and order as \code{coef}. Rows and columns of fixed
parameters are \code{NA}.
\item[\code{mask}] Logical vector with the same length and order as \code{coef},
\code{TRUE} for fixed parameters.
\item[\code{loglik}] Log-likelihood of the fitted model.
\item[\code{aic}] Akaike Information Criterion.
\item[\code{bic}] Bayesian Information Criterion.
\item[\code{aicc}] Corrected Akaike Information Criterion.
\item[\code{sigma2}] Estimated innovation variance.
\item[\code{x}] The original time series.
\item[\code{fitted}] One-step-ahead fitted values, on the scale of the original series.
\item[\code{residuals}] Model residuals. Together with \code{fitted}, the first
\code{d + sum(D * S)} values are \code{NA}, since differencing leaves no
observations to predict for them.
\item[\code{method}] Character string describing the fitted model, e.g.
\code{"mSARIMA(1,0,0)(1,0,0)[4]"}.
\item[\code{series}] Name of the series.
\item[\code{nobs}] Number of observations in the original series.
\item[\code{n.used}] Number of observations used in estimation (after differencing).
\item[\code{call}] The matched call.
\item[\code{convergence}] Convergence code from \code{optim} (0 indicates successful convergence).
\end{ldescription}
\end{Value}
%
\begin{Note}
\begin{itemize}

\item{} The series must have at least 10 observations, both before and after differencing.
\item{} Seasonal periods must be at least 2 and at most half the length of the series.
\item{} If \code{y} is a \code{ts} or \code{msts} object, \code{x}, \code{fitted} and
\code{residuals} keep the same time series structure. If \code{y} is a plain numeric
vector, they are returned as \code{ts} with one seasonal period, or as \code{msts}
with several.

\end{itemize}

\end{Note}
%
\begin{Examples}
\begin{ExampleCode}
# Example 1: ARIMA with two seasonal components
set.seed(123)
y <- msarima.sim(nsim = 200,
                 order = c(1, 0, 0),
                 seasonal = list(c(1, 0, 0), c(0, 0, 1)),
                 seasonal.period = c(4, 11),
                 ar = c(0.5, 0.5), ma = 0.3, sigma2 = 1)

fit <- msarima(y, order = c(1, 0, 0),
               seasonal = list(c(1, 0, 0), c(0, 0, 1)),
               seasonal.period = c(4, 11))
print(fit)
summary(fit)

# Example 2: msts input, seasonal periods detected automatically
y_msts <- msts(y, seasonal.periods = c(4, 11))
fit_msts <- msarima(y_msts, order = c(1, 0, 0),
                    seasonal = list(c(1, 0, 0), c(0, 0, 1)))

# Example 3: differencing
fit_diff <- msarima(cumsum(y), order = c(1, 1, 0),
                    seasonal = list(c(0, 0, 1)),
                    seasonal.period = 4)
head(residuals(fit_diff))  # the first residual is NA

# Example 4: fixing parameters
# Note: the length of 'fixed' must match AR + SAR + MA + SMA + sigma
fit_fixed <- msarima(y, order = c(1, 0, 0),
                     seasonal = list(c(1, 0, 0), c(0, 0, 1)),
                     seasonal.period = c(4, 11),
                     fixed = c(0.5, NA, NA, NA))
AIC(fit); AIC(fit_fixed)

\end{ExampleCode}
\end{Examples}
\HeaderA{msarima-methods}{Methods for msarima Objects}{msarima.Rdash.methods}
\aliasA{AIC.msarima}{msarima-methods}{AIC.msarima}
\aliasA{BIC.msarima}{msarima-methods}{BIC.msarima}
\aliasA{coef.msarima}{msarima-methods}{coef.msarima}
\aliasA{fitted.msarima}{msarima-methods}{fitted.msarima}
\aliasA{logLik.msarima}{msarima-methods}{logLik.msarima}
\aliasA{print.msarima}{msarima-methods}{print.msarima}
\aliasA{print.summary\_msarima}{msarima-methods}{print.summary.Rul.msarima}
\aliasA{residuals.msarima}{msarima-methods}{residuals.msarima}
\aliasA{summary.msarima}{msarima-methods}{summary.msarima}
%
\begin{Description}
Methods for extracting information from and displaying fitted multi-seasonal ARIMA models.
\end{Description}
%
\begin{Usage}
\begin{verbatim}
## S3 method for class 'msarima'
print(x, ...)

## S3 method for class 'msarima'
summary(object, ...)

## S3 method for class 'summary_msarima'
print(x, digits = 4, signif.stars = TRUE, ...)

## S3 method for class 'msarima'
coef(object, ...)

## S3 method for class 'msarima'
fitted(object, ...)

## S3 method for class 'msarima'
residuals(object, ...)

## S3 method for class 'msarima'
logLik(object, ...)

## S3 method for class 'msarima'
AIC(object, ...)

## S3 method for class 'msarima'
BIC(object, ...)
\end{verbatim}
\end{Usage}
%
\begin{Arguments}
\begin{ldescription}
\item[\code{x}, \code{object}] An object of class \code{"msarima"}.

\item[\code{...}] Additional arguments (currently not used).

\item[\code{digits}] Number of digits to display. Default is 4.

\item[\code{signif.stars}] Logical; if \code{TRUE}, significance stars are printed. Default is \code{TRUE}.
\end{ldescription}
\end{Arguments}
%
\begin{Details}
\begin{description}

\item[\code{print}] Prints a concise summary of the fitted model including coefficients
and information criteria.
\item[\code{summary}] Produces a detailed summary including coefficient estimates,
standard errors, z-values, and p-values.
\item[\code{coef}] Extracts the estimated model coefficients.
\item[\code{fitted}] Extracts the fitted values.
\item[\code{residuals}] Extracts the model residuals.
\item[\code{logLik}] Extracts the log-likelihood value.
\item[\code{AIC}] Calculates the Akaike Information Criterion.
\item[\code{BIC}] Calculates the Bayesian Information Criterion.

\end{description}

\end{Details}
%
\begin{Value}
\begin{description}

\item[\code{print}] Returns the input object invisibly.
\item[\code{summary}] Returns an object of class \code{"summary\_msarima"} containing:
\begin{itemize}

\item{} \code{call}: The matched call.
\item{} \code{series}: Name of the time series.
\item{} \code{order}: Non-seasonal ARIMA order.
\item{} \code{seasonal}: List of seasonal orders (P, D, Q), one vector per
seasonal component.
\item{} \code{seasonal.period}: Vector of seasonal periods.
\item{} \code{coefficients}: Matrix with estimates, standard errors, z-values, and p-values.
\item{} \code{sigma2}: Estimated innovation variance.
\item{} \code{loglik}: Log-likelihood of the fitted model.
\item{} \code{aic}: Akaike Information Criterion.
\item{} \code{bic}: Bayesian Information Criterion.
\item{} \code{aicc}: Corrected Akaike Information Criterion.
\item{} \code{nobs}: Number of observations in the original series.
\item{} \code{n.used}: Number of observations used in estimation.
\item{} \code{method}: Character string describing the fitted model, e.g.
\code{"mSARIMA(1,0,0)(1,0,0)[4]"}.
\item{} \code{convergence}: Convergence code from optimization.

\end{itemize}


\item[\code{coef}] A named numeric vector of coefficients.
\item[\code{fitted}] A time series object of fitted values.
\item[\code{residuals}] A time series object of residuals.
\item[\code{logLik}] An object of class \code{"logLik"}.
\item[\code{AIC}] The AIC value.
\item[\code{BIC}] The BIC value.

\end{description}

\end{Value}
%
\begin{SeeAlso}
\code{\LinkA{msarima}{msarima}}
\end{SeeAlso}
%
\begin{Examples}
\begin{ExampleCode}
# Fit a model
set.seed(123)
y <- msarima.sim(nsim = 200,
                 order = c(1, 0, 0),
                 seasonal = list(c(1, 0, 0), c(0, 0, 1)),
                 seasonal.period = c(4, 11),
                 ar = c(0.5, 0.5), ma = 0.3, sigma2 = 1)

fit <- msarima(y, order = c(1, 0, 0),
               seasonal = list(c(1, 0, 0), c(0, 0, 1)),
               seasonal.period = c(4, 11))

# Use methods
print(fit)
summary(fit)
coef(fit)
head(fitted(fit))
head(residuals(fit))
logLik(fit)
AIC(fit)
BIC(fit)

\end{ExampleCode}
\end{Examples}
\HeaderA{msarima.sim}{Simulate from a Multi-Seasonal ARIMA Model (Direct Interface)}{msarima.sim}
%
\begin{Description}
Simulates a time series from a multi-seasonal ARIMA model with specified parameters.
This function provides a direct interface for simulation without requiring a fitted model object.
\end{Description}
%
\begin{Usage}
\begin{verbatim}
msarima.sim(
  nsim,
  order = c(0, 0, 0),
  seasonal = list(),
  seasonal.period = NULL,
  ar = NULL,
  ma = NULL,
  sigma2 = 1,
  burn = NULL,
  intercept = 0
)
\end{verbatim}
\end{Usage}
%
\begin{Arguments}
\begin{ldescription}
\item[\code{nsim}] Length of the simulated series (required).

\item[\code{order}] A numeric vector of length 3 specifying the non-seasonal part of the
ARIMA model: the AR order (p), the degree of differencing (d), and the MA order (q).
Default is \code{c(0, 0, 0)}.

\item[\code{seasonal}] A list of numeric vectors, each of length 3, specifying the seasonal
components. Each vector gives the seasonal AR order (P), the seasonal differencing
order (D), and the seasonal MA order (Q) of one component. Use \code{list()} for a
model without seasonal components. Default is \code{list()}.

\item[\code{seasonal.period}] A numeric vector of seasonal periods, one for each element of
\code{seasonal}. Periods can be given in any order but must be distinct.
Default is \code{NULL}.

\item[\code{ar}] Numeric vector of AR coefficients, non-seasonal first and then one block per
seasonal component. Its length must match the total AR order implied by \code{order}
and \code{seasonal}. Default is \code{NULL}.

\item[\code{ma}] Numeric vector of MA coefficients, ordered as \code{ar}. Its length must match
the total MA order implied by \code{order} and \code{seasonal}. Default is \code{NULL}.

\item[\code{sigma2}] Innovation variance. Default is 1.

\item[\code{burn}] Number of burn-in observations to discard. If \code{NULL}, it is set to
\code{max(100, 2 * kk)}, where \code{kk} is the largest lag in the model.
Default is \code{NULL}.

\item[\code{intercept}] Constant term of the model equation, that is \eqn{\mu}{} in
\eqn{y_t = \mu + \sum \phi_i y_{t-i} + \sum \theta_j e_{t-j} + e_t}{}. This is the
intercept, not the mean of the process: for a stationary model the mean is
\code{intercept / prod(1 - ar)}. Default is 0.
\end{ldescription}
\end{Arguments}
%
\begin{Value}
A time series object of length \code{nsim}: \code{ts} with one seasonal period,
\code{msts} with several.
\end{Value}
%
\begin{SeeAlso}
\code{\LinkA{msarima}{msarima}}, \code{\LinkA{simulate.msarima}{simulate.msarima}}
\end{SeeAlso}
%
\begin{Examples}
\begin{ExampleCode}
# Simulate data
set.seed(123)
y <- msarima.sim(nsim = 200,
                 order = c(1, 0, 0),
                 seasonal = list(c(1, 0, 0), c(0, 0, 1)),
                 seasonal.period = c(4, 11),
                 ar = c(0.5, 0.5), ma = 0.3, sigma2 = 1)

\end{ExampleCode}
\end{Examples}
\HeaderA{msportmanteau}{Multi-Seasonal Portmanteau Tests}{msportmanteau}
%
\begin{Description}
Multi-seasonal generalisation of the Ljung-Box and Monti portmanteau tests:
residual autocorrelation is checked at low-order and seasonal lags at once,
for an arbitrary number of seasonal periods.
\end{Description}
%
\begin{Usage}
\begin{verbatim}
msportmanteau(
  x,
  type = c("Ljung-Box", "Monti"),
  season = NULL,
  nlag = 5,
  lags_short = 0,
  npar = NULL
)
\end{verbatim}
\end{Usage}
%
\begin{Arguments}
\begin{ldescription}
\item[\code{x}] A numeric vector, a \code{ts}/\code{msts} object, or a fitted \LinkA{msarima}{msarima} model
(in which case the model residuals are used).

\item[\code{type}] Either \code{"Ljung-Box"} (default), which builds the statistic from
the autocorrelations, or \code{"Monti"}, which uses the partial
autocorrelations. Partial matching is allowed.

\item[\code{season}] Numeric vector of seasonal periods to test, of any length. If
\code{NULL} (default) the periods are taken from the \code{msts}/\code{ts} attributes of
\code{x}, or from the fitted model.

\item[\code{nlag}] Numeric vector giving the number of seasonal multiples of each
period to test, i.e. lags \code{season * (1:nlag)}. A single value is recycled across
\code{season}; otherwise give one value per period, e.g. \code{c(5, 3)}. Default \code{5}.

\item[\code{lags\_short}] Number of low-order, non-seasonal lags \code{1:lags\_short} added
to the seasonal ones. Default \code{0} (none).

\item[\code{npar}] Number of estimated model parameters. Default \code{0} for a raw
series; for a fitted \LinkA{msarima}{msarima} model it defaults to the number of estimated
(non-fixed) AR, seasonal AR, MA and seasonal MA coefficients.
\end{ldescription}
\end{Arguments}
%
\begin{Details}
The lag set tested is \code{sort(unique(c(1:lags\_short, outer(season, 1:nlag))))},
the same construction as in the original two-season versions of these
tests, generalised to any number of periods.

As in \LinkA{sacf}{sacf}, leading and trailing \code{NA}s (such as those produced by
differencing in \LinkA{msarima}{msarima}) are dropped (\LinkA{stats::na.contiguous}{stats::na.contiguous}) before
computing the statistic, and the number of observations used is the length
of the series that remains.
\end{Details}
%
\begin{Value}
An object of class \code{"htest"}, with:
\begin{ldescription}
\item[\code{statistic}] the test statistic ("X-squared").
\item[\code{parameter}] the degrees of freedom ("df"), the number of lags tested
minus \code{npar}.
\item[\code{p.value}] the p-value, from the upper tail of a chi-squared
distribution with \code{parameter} degrees of freedom.
\item[\code{method}] a character string describing the test.
\item[\code{data.name}] a character string naming the data.
\item[\code{lag}] the lags actually tested.
\end{ldescription}
\end{Value}
%
\begin{References}
Grigoletto, M. and Lisi, F. (2026). Portmanteau goodness-of-fit tests in the
presence of simple, mixed, and multiple seasonal autocorrelation. \emph{arXiv}.
\Rhref{https://doi.org/10.48550/arXiv.2609.32381}{doi:10.48550\slash{}arXiv.2609.32381}

Ljung, G. M. and Box, G. E. P. (1978). On a measure of lack of fit in time
series models. \emph{Biometrika}, 65(2), 297-303.

Monti, A. C. (1994). A proposal for a residual autocorrelation test in
linear models. \emph{Biometrika}, 81(4), 776-780.
\end{References}
%
\begin{SeeAlso}
\LinkA{sacf}{sacf}, \LinkA{msarima}{msarima}, \LinkA{stats::Box.test}{stats::Box.test}
\end{SeeAlso}
%
\begin{Examples}
\begin{ExampleCode}
fit <- msarima(MAD_m40, order = c(1, 0, 2),
               seasonal = list(c(0, 1, 1), c(0, 1, 1)))

# periods 24 and 168 taken from the msts object
msportmanteau(MAD_m40, nlag = 3)
# residual test, periods and npar from the model
msportmanteau(fit, type = "Monti", nlag = 3)

\end{ExampleCode}
\end{Examples}
\HeaderA{sacf}{Autocorrelation Plots for Multiple Seasonal Cycles}{sacf}
\aliasA{ggsacf}{sacf}{ggsacf}
\aliasA{ggspacf}{sacf}{ggspacf}
\aliasA{spacf}{sacf}{spacf}
%
\begin{Description}
Sample autocorrelation (\code{sacf}) and partial autocorrelation (\code{spacf}) functions
that highlight the lags belonging to one or more seasonal cycles, so that the
seasonal structure of a multi-seasonal series is visible at a glance.
\code{ggsacf()} and \code{ggspacf()} are the \pkg{ggplot2} equivalents.
\end{Description}
%
\begin{Usage}
\begin{verbatim}
sacf(
  x,
  season = NULL,
  lag.max = NULL,
  ci = 0.95,
  col = NULL,
  main = NULL,
  ylab = "ACF",
  plot = TRUE,
  ...
)

spacf(x, ..., ylab = "PACF")

ggsacf(
  x,
  season = NULL,
  lag.max = NULL,
  ci = 0.95,
  main = NULL,
  ylab = "ACF",
  ...
)

ggspacf(x, ..., ylab = "PACF")
\end{verbatim}
\end{Usage}
%
\begin{Arguments}
\begin{ldescription}
\item[\code{x}] A numeric vector, a \code{ts}/\code{msts} object, or a fitted \LinkA{msarima}{msarima} model
(in which case the model residuals are used).

\item[\code{season}] Numeric vector of seasonal periods to highlight, of any length.
If \code{NULL} (default) the periods are taken from the \code{msts}/\code{ts} attributes of
\code{x}, or from the fitted model.

\item[\code{lag.max}] Maximum lag. Defaults to \code{3 * max(season)}, or to the
\LinkA{stats::acf}{stats::acf} default when no seasonal period is available.

\item[\code{ci}] Coverage of the confidence bands, as a probability. Default \code{0.95}.

\item[\code{col}] Colours used for the seasonal lags, one per element of \code{season}
(recycled). Defaults to a colour-blind safe palette (Okabe and Ito, 2008),
the same one used by \code{ggsacf()} and \code{ggspacf()}.

\item[\code{main}, \code{ylab}] Plot title and y-axis label.

\item[\code{plot}] If \code{TRUE} (default) the correlogram is drawn.

\item[\code{...}] Further arguments passed to \LinkA{stats::acf}{stats::acf}, and from \code{spacf()} and
\code{ggspacf()} to \code{sacf()} and \code{ggsacf()}.
\end{ldescription}
\end{Arguments}
%
\begin{Details}
Lags that are multiples of more than one period (lag 168 with
\code{season = c(24, 168)}, say) are attributed to the longest matching period.

The series is treated as a plain numeric vector, so lags are always reported in
number of observations and not in units of time. Leading and trailing \code{NA}s,
such as those produced by differencing in \LinkA{msarima}{msarima}, are dropped
(\LinkA{stats::na.contiguous}{stats::na.contiguous}).
\end{Details}
%
\begin{Value}
\code{sacf()} and \code{spacf()} invisibly return the \LinkA{stats::acf}{stats::acf} object, with
the seasonal periods in the \code{"season"} attribute and, for each lag, the index
of the period it belongs to in \code{"season.group"}. \code{ggsacf()} and \code{ggspacf()}
return a \code{ggplot} object.
\end{Value}
%
\begin{References}
Okabe, M. and Ito, K. (2008). Color Universal Design (CUD): How to make figures
and presentations that are friendly to colorblind people.
\url{https://jfly.uni-koeln.de/color/}
\end{References}
%
\begin{SeeAlso}
\LinkA{stats::acf}{stats::acf}, \LinkA{msarima}{msarima}
\end{SeeAlso}
%
\begin{Examples}
\begin{ExampleCode}
fit <- msarima(MAD_m40, order = c(1, 0, 2),
               seasonal = list(c(0, 1, 1), c(0, 1, 1)))

sacf(MAD_m40, lag.max = 200)   # periods 24 and 168 taken from the msts object
spacf(fit)                     # residual pacf, periods taken from the model

\end{ExampleCode}
\end{Examples}
\HeaderA{simulate.msarima}{Simulate from a Multi-Seasonal ARIMA Model}{simulate.msarima}
%
\begin{Description}
Returns a simulated time series from a fitted ARIMA model with one or more seasonal
components.
\end{Description}
%
\begin{Usage}
\begin{verbatim}
## S3 method for class 'msarima'
simulate(object = NULL, nsim = 100, seed = NULL, burn = NULL, ...)
\end{verbatim}
\end{Usage}
%
\begin{Arguments}
\begin{ldescription}
\item[\code{object}] An object of class \code{"msarima"}.

\item[\code{nsim}] Length of the simulated series. Default is 100.

\item[\code{seed}] Either \code{NULL} or an integer passed to \code{\LinkA{set.seed}{set.seed}}
before simulating, to make the result reproducible. Default is \code{NULL}.

\item[\code{burn}] Number of burn-in observations to discard. If \code{NULL}, it is set by
\code{\LinkA{msarima.sim}{msarima.sim}} to \code{max(100, 2 * kk)}, where \code{kk} is the
largest lag in the model. Default is \code{NULL}.

\item[\code{...}] Additional arguments (currently not used).
\end{ldescription}
\end{Arguments}
%
\begin{Details}
The model specification and the estimated parameters, including the constant term when
the model was fitted with \code{include.constant = TRUE}, are extracted from the fitted
object and passed to \code{\LinkA{msarima.sim}{msarima.sim}}.
\end{Details}
%
\begin{Value}
A time series object (\code{ts} or \code{msts}) of length \code{nsim}, with the
seasonal periods of the fitted model.
\end{Value}
%
\begin{SeeAlso}
\code{\LinkA{msarima}{msarima}}, \code{\LinkA{msarima.sim}{msarima.sim}}
\end{SeeAlso}
%
\begin{Examples}
\begin{ExampleCode}
# Simulate data
set.seed(123)
y <- msarima.sim(nsim = 200,
                 order = c(1, 0, 0),
                 seasonal = list(c(1, 0, 0), c(0, 0, 1)),
                 seasonal.period = c(4, 11),
                 ar = c(0.5, 0.5), ma = 0.3, sigma2 = 1)

# Fit model
fit <- msarima(y, order = c(1, 0, 0),
               seasonal = list(c(1, 0, 0), c(0, 0, 1)),
               seasonal.period = c(4, 11))

# Simulate from fitted model
y_sim <- simulate(fit, nsim = 100)

\end{ExampleCode}
\end{Examples}
\printindex{}
\end{document}
